In recent years, machine learning-based \cite{chen2020zero,chen2021one,chen2022representative,chen2021ct,wen2020noises} medical diagnosis systems have been introduced by many institutions. Although these systems achieve high accuracy in predicting medical conditions, bias has been found in dermatological disease datasets as shown in \cite{groh2021evaluating,kinyanjui2020fairness,tschandl2018ham10000}. This bias can arise when there is an imbalance in the number of images representing different skin tones, which can lead to inaccurate predictions and misdiagnosis due to biases towards certain skin tones. The discriminatory nature of these models can have a negative impact on society by limiting access to healthcare resources for different sensitive groups, such as those based on race or gender.

% 要跟老師討論，因為這邊的敘述是我們之前直接抄fairprune的 加上reviewer有疑問 所以要重寫一下比較妥

% remove sensitive information \cite{wang2020towards} and
Several methods are proposed to alleviate the bias in machine learning models, including pre-processing, in-processing, and post-processing strategies. Pre-processing strategies adjust training data before training \cite{lu2020gender,ngxande2020bias} or assign different weights to different data samples to suppress the sensitive information during training \cite{kamiran2012data}. In-processing modifies the model architecture, training strategy, and loss function to achieve fairness, such as adversarial training \cite{alvi2018turning,zhang2018mitigating} or regularization-based \cite{jung2021fair,quadrianto2019discovering} methods. Recently, pruning \cite{wu2022fairprune} techniques have also been used to achieve fairness in dermatological disease diagnosis. However, these methods may decrease accuracy for both groups and do not guarantee explicit protection for the unprivileged group when enforcing fairness constraints. Post-processing techniques enhance fairness by adjusting the model's output distribution. This calibration is done by taking the model's prediction, and the sensitive attribute as inputs \cite{du2020fairness,hardt2016equality,zhao2017men}. However, pre-processing and post-processing methods have limitations that are not applicable to dermatological disease diagnostic tasks since they need extra sensitive information during the training time.
 
% Recently, many institutions have introduced machine learning-based medical diagnosis systems. Although the machine learning model achieves accurate prediction, there exists some algorithmic bias on two dermatological datasets ISIC2018  \cite{liu2018large} and Fitzpatrick-17k \cite{Bengio+chapter2007} datasets shown in \cite{groh2021evaluating,kinyanjui2020fairness}. Furthermore, the discriminative nature of the machine learning model will harm the opportunity of different protected groups (e.g., race, gender) to access healthcare resources and thus tear society apart.

% Several methods are proposed to ameliorate the bias in machine learning models. Many of them \cite{wang2022fairness, zhang2018mitigating, kim2019learning, ngxande2020bias} adopted adversarial training to eliminate bias by training the network to learn a classifier while disabling the adversary's ability to categorize the sensitive attribute. Disentanglement representation \cite{creager2019flexibly} is another mainstream to achieve fairness. It forces the latent vector of the sensitive group to be independent of that of the target group and thus reaches fairness.
% Several methods are proposed to ameliorate the bias in machine learning models. Many of them \cite{wang2022fairness} \cite{Wang2020TowardsFI} \cite{Zhang2018MitigatingUB} \cite{kim2019learning} \cite{Ngxande2020BiasRI} adopted adversarial training to eliminate bias by training the network to learn a classifier while disabling the adversary's ability to categorize the sensitive attribute. Disentanglement representation \cite{creager2019flexibly} is another mainstream to achieve fairness. It forces the latent vector of the sensitive group to be independent of that of the target group and thus reach fairness.

In this paper, we observe that although features from a deep layer of a neural network are discriminative for target groups (i.e., different dermatological diseases), they cause fairness conditions to deteriorate, and we will demonstrate this observation by analyzing the entanglement degree regarding sensitive information with the soft nearest neighbor loss \cite{frosst2019analyzing} of image features in Section \ref{sec:motivation}. This finding is similar to ``overthinking'' phenomenon in neural networks \cite{kaya2019shallow}, where accuracy decreases as the features come from deeper in a neural network and motivate us to use a multi-exit network \cite{kaya2019shallow,teerapittayanon2016branchynet} to address the fairness issue.

% where accuracy decreases as the features come from deeper in a neural network

Through extensive experiments, we demonstrate that our proposed multi-exit convolutional neural network (ME-CNN) can achieve fairness without using sensitive attributes (unawareness) in the training process, which is suitable for dermatological disease diagnosis because the sensitive attributes information exists privacy issues and is not always available. We compare our approach to the current state-of-the-art method proposed in \cite{wu2022fairprune} and find that the ME-CNN can achieve similar levels of fairness. To further improve fairness, we designed a new framework for fair multi-exit. With the fairness constraint and the early exit strategy at the inference stage, a sufficient discriminative basis can be obtained based on low-level features when classifying easier samples. This contributes to selecting a more optimal prediction regarding the trade-off between accuracy and fairness for each test instance.

% This finding reminds us of the ``overthinking'' phenomenon in deep neural networks \cite{kaya2019shallow}, where accuracy decreases as the features come from deeper in a neural network. This problem is successfully addressed through multi-exit neural networks by introducing multiple internal classifiers and treating high-confidence results from these internal classifiers as the final result. Therefore, a similar approach can be used to address the issue concerning fairness. 

% In the proposed multi-exit framework for fairness, we leverage the early exits strategy at the inference stage; a sufficient discriminative basis can be obtained based on low-level features when classifying easier samples. This contributes to selecting a more optimal prediction regarding the trade-off between accuracy and fairness for each test instance. To validate the effectiveness of our framework on medical images, we perform the dermatological image classification on ISIC2019  \cite{liu2018large} and Fitzpatrick-17k \cite{Bengio+chapter2007} datasets. Experiments show that the proposed method significantly improves the results from the baseline and different state-of-the-art in terms of the trade-off between classification accuracy and fairness.

The main contributions of the proposed method are as follows: 

\begin{itemize}
    % \item We observe that the features from a deep layer of a neural network are highly discriminative yet cause fairness to deteriorate.
    \item Our quantitative analysis shows that the features from a deep layer of a neural network are highly discriminative yet cause fairness to deteriorate.
    \item We propose a fairness through unawareness framework and use multi-exit training to improve fairness in dermatological disease classification.
    \item We demonstrate the extensibility of our framework, which can be applied to various state-of-the-art models and achieve further improvement. Through extensive experiments, we show that our approach can improve fairness while keeping competitive accuracy on both the dermatological disease dataset, ISIC 2019 \cite{combalia2019bcn20000}, and Fitzpatrick-17k datasets \cite{groh2021evaluating}. 
    % \item Through extensive experiments on different settings and comparisons, our framework achieves the best trade-off performances between top-1 accuracy and fairness on ISIC2018 and Fitzpatrick-17k datasets.

\end{itemize}

% \begin{itemize}
%     \item Our quantitative analysis shows that the features from a deep layer of a neural network are highly discriminative yet cause the fairness to deteriorate.
%     \item We explore the use of a multi-exit training framework to deal with the fairness issue.
%     \item We introduce a simple fairness through unawareness framework for dermatological image classification tasks with high extensibility that could apply to different baselines and state-of-the-art and achieve further improvement.
%     % \item Through extensive experiments on different settings and comparisons, our framework achieves the best trade-off performances between top-1 accuracy and fairness on ISIC2018 and Fitzpatrick-17k datasets.

% \end{itemize}






% Recently, many institutions have introduced machine learning-based medical diagnosis systems. Although the machine learning model achieves accurate prediction, there exists some algorithmic bias on two dermatological datasets ISIC2018  \cite{liu2018large} and Fitzpatrick-17k \cite{Bengio+chapter2007} datasets shown in \cite{groh2021evaluating,kinyanjui2020fairness}. The discriminative nature of the machine learning model will harm the opportunity of different protected groups (e.g., race, gender) to access healthcare resources and thus tear society apart.

% % The discriminative nature of the machine learning model will harm the opportunity of different races, religions, and genders and thus tear society apart.

% % Several methods are proposed to ameliorate the bias in machine learning models. Many of them \cite{wang2022fairness, zhang2018mitigating, kim2019learning, ngxande2020bias} adopted adversarial training to eliminate bias by training the network to learn a classifier while disabling the adversary's ability to categorize the sensitive attribute. Disentanglement representation \cite{creager2019flexibly} is another mainstream to achieve fairness. It forces the latent vector of the sensitive group to be independent of that of the target group and thus reaches fairness.
% % Several methods are proposed to ameliorate the bias in machine learning models. Many of them \cite{wang2022fairness} \cite{Wang2020TowardsFI} \cite{Zhang2018MitigatingUB} \cite{kim2019learning} \cite{Ngxande2020BiasRI} adopted adversarial training to eliminate bias by training the network to learn a classifier while disabling the adversary's ability to categorize the sensitive attribute. Disentanglement representation \cite{creager2019flexibly} is another mainstream to achieve fairness. It forces the latent vector of the sensitive group to be independent of that of the target group and thus reach fairness.

% In this paper, we observe that although features from a deep layer of a neural network are discriminative for target groups (i.e., different dermatological diseases), they cause fairness conditions to deteriorate, and we will demonstrate this observation by analyzing the entanglement degree regarding sensitive information with the soft nearest neighbor loss \cite{frosst2019analyzing} of image features in Section \ref{sec:motivation}. 

% This finding reminds us of the ``overthinking'' phenomenon in deep neural networks \cite{kaya2019shallow}, where accuracy decreases as the features come from deeper in a neural network. This problem is successfully addressed through multi-exit neural networks by introducing multiple internal classifiers and treating high-confidence results from these internal classifiers as the final result. We conjecture that a similar approach can be used to address the issue concerning fairness. 

% In the proposed multi-exit framework for fairness, both accuracy and fairness constraints are included when training every internal classifier to keep the fairness and accuracy from shallow to deep. Specifically, with early exits at the inference stage, a sufficient discriminative basis can be obtained based on low-level features when classifying easier samples. This contributes to selecting a more optimal prediction for each test instance regarding the trade-off between accuracy and fairness.

% To validate the effectiveness of our framework on medical images, we perform the dermatological image classification on ISIC2018  \cite{liu2018large} and Fitzpatrick-17k \cite{Bengio+chapter2007} datasets. Experiments show that the proposed method significantly improves the results from the baseline and different state-of-the-art in terms of the trade-off between classification accuracy and fairness.

% The main contributions of the proposed method are as follows: 
% \begin{itemize}
%     \item Our quantitative analysis shows that the features from a deep layer of a neural network are highly discriminative yet cause the fairness to deteriorate.
%     \item We explore the use of multi-exit training frameworks to deal with the fairness issue.
%     \item We introduce a simple fairness through unawareness framework for dermatological image classification tasks with high extensibility that could apply to different baselines and state-of-the-art and achieve further improvement.
%     % \item Through extensive experiments on different settings and comparisons, our framework achieves the best trade-off performances between top-1 accuracy and fairness on ISIC2018 and Fitzpatrick-17k datasets.

% \end{itemize}



